\section{Methods}\label{methods}

\subsection{Deep Neural Networks}

A Deep Neural Network (DNN) is a foundational machine learning architecture modelled on the human brain. Featuring multiple ``hidden" layers of interconnected nodes, it processes complex, unstructured data like text, images, and audio to autonomously discover patterns, extract features, and make accurate predictions or decisions. \cite{deeplearningreview} \cite{Goodfellow-et-al-2016}

\begin{figure}[!htbp]
\centering
\includegraphics[scale=0.4]{deepneuralnetwork}
\caption{A Deep Neural Network}
\end{figure}

\begin{itemize}

\item DNNs are organized into three distinct types of layers, each playing a specific role in data transformation:

\begin{enumerate}
    \item \textbf{Input Layer:} Receives the raw data (e.g., individual pixels of an image or numerical values in a dataset) and passes it to the network.
    
    \item \textbf{Hidden Layers:} These are the multiple layers between the input and output layers. They perform nonlinear transformations to extract progressively complex features.
    
    \item \textbf{Output Layer:} Produces the final prediction, classification, or decision.
\end{enumerate}

\item The learning process in a DNN generally consists of the following steps:

\begin{enumerate}
    \item \textbf{Forward Propagation:} Data passes through the network where each connection applies weights and biases to compute activations.
    
    \item \textbf{Loss Calculation:} The predicted output is compared with the actual target to compute the error or loss.
    
    \item \textbf{Backpropagation:} The network updates its weights by propagating the error backward in order to minimize the loss function.
\end{enumerate}

\end{itemize}

Some common DNN architectures are Convolutional Neural Networks, Recurrent Neural Networks etc.



\subsection{Residual Neural Networks}

As neural networks become deeper, they suffer from problems such as :

\begin{enumerate}
    \item \textbf{Vanishing Gradient Problem} : During backpropagation, gradients can become extremely small, making earlier layers learn very slowly.

    \item \textbf{Degradation Problem} : Adding more layers sometimes leads to worse training accuracy, even though deeper networks should theoretically perform better.

    \item \textbf{Difficulty in Training Very Deep Networks} : Traditional deep networks become unstable and hard to optimize beyond a certain depth.
\end{enumerate}

If often turns out that learning the residual is easier than learning the direct mapping. A Residual Neural Network (ResNet) \cite{he2015deepresiduallearningimage} is a neural network architecture that introduces skip connections to learn residual mappings instead of directly learning complex transformations.

\begin{figure}[!htbp]
\centering
\includegraphics[scale=0.5]{ResNet_Structure}
\caption{ResNet Structure}
\end{figure}

\begin{itemize}
\item The mapping of $f(\bs{x}_t; \bs{\theta})$ can be realised by feed-forward neural networks
\item $\bs{x}_t$ is added to the output of the stacked layers
\item The hidden state $\bs{x}_t$ of the $t-$th layer of the network can be described by the equation 
\begin{equation}\label{ResNet_hidden_state_representation}
\bs{x}_{t+1} = \bs{x}_t + f(\bs{x}_t; \bs{\theta}); \qquad t = \{0, 1, 2, \ldots, T\}.
\end{equation}
\end{itemize}


\subsection{ResNet to Ordinary Differential Equations}

A growing amount of research arises in finding the relationship between neural network structure and discretized solution of numerical differential equations. For example, as neural network maps the input $x$ to the output $y$ through a series of hidden layers, the hidden representation can be regarded as the state of a dynamic system. In particular, ResNet can be regarded as a form of discretizee Euler
solution of ODEs.

Rewrite (\ref{ResNet_hidden_state_representation}) as

$$\boldsymbol{x}(t+\Delta t) = \boldsymbol{x}(t)+\Delta t\, f(\boldsymbol{x}(t); \boldsymbol{\theta}).$$

If we \textit{add more layers} and \textit{take small steps} \textit{i.e.} $\Delta t \rightarrow 0$, then the
continuous dynamics of hidden units can be shown as follows \cite{chen2019neuralordinarydifferentialequations}

\begin{align}
\dfrac{\bs{x}(t + \Delta t) - \bs{x}(t)}{\Delta t} &= f(\bs{x}(t); \bs{\theta}) \nonumber \\[0.5em]
\Rightarrow \dfrac{d \bs{x}(t)}{dt} &= f(\bs{x}(t); \bs{\theta}).
\end{align}

Thus we see that ResNet can indeed match the \textit{Euler discretization} of ordinary differential equations.


\subsection{Neural Ordinary Differential Equations}

A Neural Ordinary Differential Equation is given by

\begin{equation}\label{NeuralODE}
\dfrac{d \bs{x}(t)}{dt} = f(\bs{x}(t); \bs{\theta})
\end{equation}

\begin{align*}
\text{where } \qquad \bs{x}(t) &: \text{ hidden representation at time }t, \\[0.2em]
f(\cdot; \bs{\theta}) &: \text{ a neural network}, \\[0.2em]
\bs{\theta} &: \text{ parameters of the neural network } f.
\end{align*}


\begin{itemize}

\item \textbf{Forward Pass}: Given \textit{initial state} $\bs{x}(0)$ and the differential equation, we compute

$$
\bs{x}(T) = \bs{x}(0) + \int \limits_0^T f(\bs{x}(t); \bs{\theta}) \; dt.
$$

$f(\bs{x}(t); \bs{\theta})$ represents \textit{how the state evolves at each instant}.

\item \textbf{Backward Pass and Learning $\bs{\theta}$}

\begin{enumerate}

\item Starting with an initial guess of $\bs{\theta}$, we calculate $\bs{x}(T)$ using an ODE solver

$$
\bs{x}(t+\Delta t)
=
\bs{x}(t)+\Delta t\, f(\bs{x}(t); \bs{\theta}).
$$

This is a \textit{Forward Pass} in practice.

\item Then compute the loss, calculate the derivatives w.r.t.\ the parameters. This is \textit{Backward Pass}.

\item Update $\bs{\theta}$ using the derivatives and repeat the process until convergence.

\end{enumerate}

\end{itemize}


\subsection{Neural Stochastic Differential Equations}

\begin{itemize}

\item \textbf{Formalization}

A Neural Stochastic Differential Equation is given by

\begin{equation}\label{NeuralSDE}
d\bs{x}(t) = f(\bs{x}(t); \bs{\theta}) \; dt + g(\bs{x}(t); \bs{\phi}) \; dW_t
\end{equation}

\begin{align*}
\text{where } \qquad \bs{x}(t) &: \text{ hidden representation at time }t, \\[0.2em]
f(\bs{x}(t); \bs{\theta}) &: \text{ \textit{drift} coefficient approximated by a neural network}, \\[0.2em]
dt &: \text{ time increment}, \\[0.2em]
g(\bs{x}(t); \bs{\phi}) &: \text{ \textit{diffusion} coefficient approximated by a neural network}, \\[0.2em]
dW_t &: \text{ random noise increment.}
\end{align*}


\item \textbf{Forward Pass}

Based on the \textit{Euler-Maruyama scheme}, a forward simulation of a SDE with Gaussian noise is approximated by

\begin{equation}\label{Euler-Maruyama_ForwardPass}
\bs{x}(t+\Delta t)
=
\bs{x}(t)
+
f(\bs{x}(t); \bs{\theta}) \; \Delta t
+
g(\bs{x}(t); \bs{\phi}) \; \sqrt{\Delta t} \; \epsilon_t.
\end{equation}

\begin{itemize}

\item Two runs with same initialization produce different paths.

\item We simulate many trajectories and take the average.

\end{itemize}

\newpage

\item \textbf{Backward Pass}

\begin{enumerate}

\item Starting at an initial state $\bs{x}(0)$, with initial approximations of $\bs{\theta}$ and $\bs{\phi}$, obtain $\bs{x}(T)$ by many small steps $\Delta t$.

\item Compute the loss and calculate the derivatives w.r.t.\ the parameters.

\item Update $\bs{\theta}$ and $\bs{\phi}$ using the derivatives and repeat until convergence.

\end{enumerate}

\end{itemize}

\subsubsection{Gaussian Noise}

\begin{itemize}
\item Depending on the nature of the random noise, there can be different classes of SDEs. 

\item The most common random noise is \textit{Gaussian Noise} under which $$dW_t \sim \mathcal{N}(0, dt).$$

\item In the forward simulation, we set $$\epsilon_t \overset{iid}{\sim} \mathcal{N}(0, 1), \text{ consequently } \sqrt{\Delta t} \; \epsilon_t \overset{iid}{\sim} \mathcal{N}(0, \Delta t).$$
\end{itemize}

\subsubsection{Non-Gaussian Noise}

Gaussian noise is not always enough to capture big jumps. Hence we introduce Non-Gaussian noise. One such noise is $\alpha$-Stable Lévy motion.\\[1em]

$\bullet$ \textbf{$\bs{\alpha}$-Stable Random Variables} : 

\begin{itemize}

\item A random variable $X$ is said to be $\alpha$-Stable if its $n$ \textit{i.i.d.} copies have the property $$\sum \limits_{i = 1}^{n} X_i \overset{d}{=} n^{1/\alpha} \; X + b_n$$ and denoted as $$X \sim S_{\alpha} (\sigma, \beta, \gamma)$$ where $\sigma :$ scale parameter, $\beta :$ skewness, $\gamma :$ location parameter.

\item $\sigma \, Z + \gamma \overset{d}{=} X$ where $Z \sim S_{\alpha}(1, \beta, 0)$.

\item $X \sim S_{\alpha}(1, 0, 0)$ is called a Standard $\alpha$-Stable random variable.

\item $\alpha \in (0, 2]$ controls how heavy tails are.

\item $\alpha = 2 \rightarrow$ Gaussian.

\item $\alpha = 1 \rightarrow$ Cauchy.

\item Smaller $\alpha \rightarrow$ heavier tails (more extreme values).
\end{itemize}


\begin{figure}[!htbp]
\centering
\includegraphics[scale=0.5]{Different Levy Distribution PDFs}
\caption{PDFs of Standard Gaussian and $\alpha-$Stable Random Variables with different $\alpha, \beta$}
\end{figure}

$\bullet$ \textbf{$\bs{\alpha}$-Stable Lévy Motions}

A symmetric $\alpha$-Stable Lévy motion $\{L_t^{\alpha}\}_{t \geq 0}$ is a stochastic process with properties

\begin{enumerate}[(i)]
\item $L_0^{\alpha} = 0$,
\vspace{0.2cm}
\item $L_t^{\alpha}$ has independent increments,
\vspace{0.2cm}
\item $L_t^{\alpha} - L_s^{\alpha} \sim S_{\alpha}\left( (t-s)^{1 / \alpha}, 0, 0 \right)$.
\end{enumerate}


\begin{itemize}
\item For a very small step $dt$, $dL_{t + dt}^{\alpha} - dL_{t}^{\alpha} = dL_t^{\alpha} \sim S_{\alpha}\left( (dt)^{1 / \alpha}, 0, 0 \right)$.

\item We can write $(dt)^{1 / \alpha} \, L_t \overset{d}{=} dL_t^{\alpha}$ where $L_t \sim S_{\alpha} (1, 0, 0)$.

\item In SDE, we set $$L_t \overset{iid}{\sim} S_{\alpha}(1, 0, 0), \text{ consequently } (\Delta t)^{1 / \alpha} \; L_t \overset{iid}{\sim} S_{\alpha}\left((\Delta t)^{1 / \alpha}, 0, 0\right).$$

\item Euler-Maruyama approximation becomes $$\boldsymbol{x}(t+\Delta t) = \boldsymbol{x}(t)+ f(\boldsymbol{x}(t); \boldsymbol{\theta}) \; \Delta t + g(\boldsymbol{x}(t); \boldsymbol{\phi}) \; (\Delta t)^{1 / \alpha} \; L_t.$$
\end{itemize}

\newpage

\begin{figure}[!htbp]
\centering
\includegraphics[scale=0.5]{Brownian Motion and Levy Motions}
\caption{Brownian Motion and Lévy Motion with different $\alpha$ and $\beta$}
\end{figure}

\subsection{LDENet}

We add Lévy motion to the right side of equation (\ref{ResNet_hidden_state_representation}) to make it the Euler discretization form of the stochastic differential equation. The continuous-time dynamics of the system (\ref{NeuralODE}) are then expressed as:
\begin{equation}\label{LDENet}
d\bs{x}(t) = f(\bs{x}(t); \bs{\theta}) \; dt + g(\bs{x}(t); \bs{\phi}) \; dL_t^{\alpha}, \; \bs{x}(0) = \bs{x}_0.
\end{equation}

Based on the Euler-Maruyama scheme, the approximation for equation (\ref{LDENet}) is

\begin{equation}\label{LDENet_Euler_Maruyama}
\boldsymbol{x}(t+\Delta t) = \boldsymbol{x}(t)+ \underbrace{f(\boldsymbol{x}(t); \boldsymbol{\theta})}_{\text{drift neural net}} \; \Delta t + \underbrace{g(\boldsymbol{x}(t); \boldsymbol{\phi})}_{\text{diffusion neural net}} \; (\Delta t)^{1 / \alpha} \; L_t.
\end{equation}

Both drift and diffusion coefficients in equation (\ref{LDENet_Euler_Maruyama}) are estimated by two different multi-layer
perceptron (MLP) neural networks which are called the drift net and diffusion net respectively. \cite{yang2022neuralnetworkstochasticdifferential}

As of now, predictive performance is of interest. So we proceed with MSELoss function.
\clearpage